\scanpage{6}

% loc: top, continuation of e
$u$. Dans ce cas, l'inclusion $E_0\hookrightarrow E$ admet un adjoint à \uline{gauche}, savoir $v_0\circ\underbrace{(u\ \text{astreint à }F_0)}_{u'}$ :

% loc: top, Hom calculation
\[
\begin{aligned}
\operatorname{Hom}_{E}(\underset{\in\operatorname{Ob}E}{Y},\underset{\in\operatorname{Ob}E_0}{X_0})
&\underset{v_0u_0\simeq\operatorname{id}}{\simeq}\operatorname{Hom}_{E}(Y,v_0u_0X_0)=\operatorname{Hom}_{E}(Y,vu_0X_0)\underset{\text{adjonction}}{\simeq}\operatorname{Hom}_{F}(uY,u_0X_0)\\
&=\operatorname{Hom}_{F_0}(u'Y,u_0X_0)\underset{v_0\,\text{pl. fid}}{\simeq}\operatorname{Hom}_{E_0}(v_0u'Y,v_0u_0X_0)\\
&\underset{v_0u_0\simeq\operatorname{id}}{=}\operatorname{Hom}_{E_0}(v_0u'Y,X_0)\qquad\text{cqfd.}
\end{aligned}
\]

% loc: f, opening
f) Conditions équivalentes :

% loc: f, items 1 and 2
1) $\operatorname{Im}v\subset E_0$ et \inserted{$1')$} $\operatorname{Im}u\subsetE_0$

2) $v\xrightarrow{\sim}vuv$ et \inserted{$2')$} $u\xleftarrow{\sim}uvu$

% loc: f, bracketed items 3 and 3 prime
[3) \deleted{$u$ se factorise $u$ se factorise en un composé} $\operatorname{Im}v\subset E_0$ \uline{et} l'inclusion $E_0\hookrightarrow E$ admet un adjoint à g

$3')$ $\operatorname{Im}u\subset F_0$ et l'inclusion $F_0\hookrightarrow F$ admet un adjoint à dr. ]

% loc: f, beside items 4 and 4 prime
% slanted
\margin{à voir si on peut simplifier}

% loc: f, item 4
4) $u$ se factorise en $E\xrightarrow{u_1}F_1\overset{\beta}{\unsure{\longrightarrow|\longleftarrow}}F$, avec $\beta$ pleinement fidèle, $u_1$ \deleted{\unsure{un épigraphe}, et $\beta$ et $u_1$ ont des adjoints à droite}\inserted{passage à catégorie des fractions [, $\beta$ et $u_1$ ayant des adjoints à dr. ]}

% loc: f, item 4 prime
$4')$ $v$ se factorise en $F\overset{v_1}{\unsure{\longrightarrow|\text{---}}} E_1\overset{\alpha}{\unsure{\longrightarrow|\text{---}}}E$, avec $\alpha$ pl. fidèle, $v_1$ passage à cat. des fractions, [ $v_1$ et $\alpha$ admettent un adj. à g. ]

% loc: lower paragraph
De \deleted{\unsure{5)|4)}} plus, la donnée d'un couple de foncteurs adjoints $(u,v)$ entre deux catégories $E,F$, satisfaisant les conditions équivalentes 1) à 4), équivaut à la donnée d'une sous-cat. str. pleine $E_0$ de $E$ telle que $E_0\overset{\alpha}{\hookrightarrow}E$ admette un adjoint à g \inserted{$\alpha'$}, d'une sous-cat. str. pleine $F_0$ de $F$ telle que $\beta:F_0\hookrightarrow F$ admette un adjoint à droite \inserted{$\beta'$}, et d'une équivalence entre $E_0$ et $F_0$

% loc: lower diagram
\diagram{\begin{tikzcd}[ampersand replacement=\&] E_0\arrow[r,shift left=.7ex,"u_0"]\&F_0\arrow[l,shift left=.7ex,"v_0"]\end{tikzcd}}

% loc: lower closing
en posant $u=\deleted{\unsure{\beta}}\beta u_0\alpha'$, $v=\alpha v_0\beta'$, les données d'adjonction étant évidentes.

% archive sheet number: 5
